%%%%%%%%%%%%%%%%%%%%%%%%%%%%%%%%%%%%%%%%%
% Guide: Creating Stata Markdown Documents
%%%%%%%%%%%%%%%%%%%%%%%%%%%%%%%%%%%%%%%%%

\documentclass[paper=a4, fontsize=11pt]{scrartcl}

\usepackage[T1]{fontenc}
\usepackage[english]{babel}
\usepackage{amsmath,amsfonts,amssymb}
\usepackage{enumerate}
\usepackage{listings}
\usepackage{xcolor}
\usepackage{mdframed}
\usepackage{graphicx}
\usepackage{booktabs}

\usepackage{sectsty}
\allsectionsfont{\normalfont\scshape}

\usepackage[hidelinks]{hyperref}
\hypersetup{colorlinks=true, citecolor=blue, urlcolor=blue, pdfborder=000}

\usepackage{fancyhdr}
\pagestyle{fancyplain}
\fancyhead{}
\fancyfoot[L]{}
\fancyfoot[C]{}
\fancyfoot[R]{\thepage}
\renewcommand{\headrulewidth}{0pt}
\renewcommand{\footrulewidth}{0pt}
\setlength{\headheight}{13.6pt}

\setlength\parindent{0pt}

% Horizontal rule
\newcommand{\horrule}[1]{\rule{\linewidth}{#1}}

% Code styling
\lstdefinestyle{stata}{
  backgroundcolor=\color{gray!8},
  basicstyle=\ttfamily\small,
  breaklines=true,
  frame=single,
  rulecolor=\color{gray!40},
  framesep=4pt,
  xleftmargin=6pt,
  xrightmargin=6pt,
}

\lstdefinestyle{shell}{
  backgroundcolor=\color{black!90},
  basicstyle=\ttfamily\small\color{white},
  breaklines=true,
  frame=single,
  rulecolor=\color{black},
  framesep=4pt,
  xleftmargin=6pt,
  xrightmargin=6pt,
}

% Tip / warning box
\newmdenv[
  backgroundcolor=blue!5,
  linecolor=blue!40,
  linewidth=1pt,
  leftmargin=0pt,
  rightmargin=0pt,
  innerleftmargin=8pt,
  innerrightmargin=8pt,
  innertopmargin=6pt,
  innerbottommargin=6pt,
]{tipbox}

\newmdenv[
  backgroundcolor=orange!8,
  linecolor=orange!60,
  linewidth=1pt,
  leftmargin=0pt,
  rightmargin=0pt,
  innerleftmargin=8pt,
  innerrightmargin=8pt,
  innertopmargin=6pt,
  innerbottommargin=6pt,
]{warnbox}

%----------------------------------------------------------------------------------------

\title{
\normalfont \normalsize
\horrule{0.5pt} \\[0.4cm]
\huge A Step-by-Step Guide to \\[0.2cm] Stata Markdown (\texttt{.stmd}) \\
\horrule{2pt} \\[0.5cm]
}

\author{Prof.\ Tzu-Ting Yang}
\date{\normalsize\today}

\begin{document}
\maketitle

\tableofcontents
\bigskip\horrule{0.4pt}\bigskip

%----------------------------------------------------------------------------------------
\section{What Is Stata Markdown?}
%----------------------------------------------------------------------------------------

A \textbf{Stata Markdown} file (\texttt{.stmd}) is a plain-text document that mixes
\textit{narrative text} (written in Markdown) with \textit{Stata code blocks}.
When you \textit{compile} it, Stata runs the code, captures all output and figures,
and a program called \textbf{Pandoc} converts the result into a polished
\textbf{HTML} (or PDF) report --- exactly like R Markdown, but for Stata.

\medskip
The workflow is:
\[
\texttt{.stmd file} \;\xrightarrow{\texttt{markstat}}\; \text{Stata runs code}
\;\xrightarrow{\texttt{Pandoc}}\; \texttt{.html} \text{ (or } \texttt{.pdf}\text{)}
\]

\bigskip

\begin{tipbox}
\textbf{Why use Stata Markdown?}
\begin{itemize}
  \item Everything --- code, output, figures, and explanation --- lives in \textit{one}
        self-contained HTML file.
  \item Results are fully \textbf{reproducible}: re-running the \texttt{.stmd} file
        always regenerates the same report.
  \item Great for homework submissions, research reports, and replication packages.
\end{itemize}
\end{tipbox}

\bigskip\horrule{0.4pt}\bigskip

%----------------------------------------------------------------------------------------
\section{Course Templates on Dropbox}
%----------------------------------------------------------------------------------------

\textbf{Goal:} produce a single \texttt{.html} file that records all your
code, output, and figures in one document --- so the reader can follow your
entire empirical workflow from data loading to final results.

\medskip
To help you get started, complete working templates are available in the
\textbf{course Dropbox folder}. Two sub-folders are relevant:

\medskip
\begin{center}
\begin{tabular}{p{4.2cm} p{9.8cm}}
\toprule
\textbf{Folder} & \textbf{Contents} \\
\midrule
\texttt{do/} &
Regular \texttt{.do} files for each method covered in class.
Use these to understand the Stata code before converting it to \texttt{.stmd}.\\[4pt]
\texttt{do/Stata\_Markdown/} &
Ready-to-compile \texttt{.stmd} templates --- one per method.
These are your \textbf{main reference} when writing your own \texttt{.stmd} file.\\
\bottomrule
\end{tabular}
\end{center}

\medskip
The available templates in \texttt{do/Stata\_Markdown/} are:

\medskip
\begin{center}
\begin{tabular}{ll}
\toprule
\textbf{File} & \textbf{Method} \\
\midrule
\texttt{regression.stmd}  & OLS regression and controls \\
\texttt{DID.stmd}         & Difference-in-Differences \\
\texttt{Stag\_DID.stmd}   & Staggered DID \\
\texttt{Syn\_DID.stmd}    & Synthetic DID \\
\texttt{fixed\_effects.stmd} & Fixed effects \\
\texttt{matching.stmd}    & Matching estimators \\
\texttt{RDD.stmd}         & Regression Discontinuity Design \\
\texttt{IV.stmd}          & Instrumental Variables \\
\texttt{SSIV.stmd}        & Shift-Share IV \\
\texttt{SCM.stmd}         & Synthetic Control Method \\
\texttt{ML.stmd}          & Machine Learning \\
\bottomrule
\end{tabular}
\end{center}

\medskip
\begin{tipbox}
\textbf{Recommended workflow:}
\begin{enumerate}
  \item Open the \texttt{.stmd} template closest to your method (e.g.\ \texttt{DID.stmd}).
  \item Read how the YAML header, markdown sections, and \texttt{```s} code blocks are
        structured.
  \item Create your own \texttt{.stmd} file following the same structure, replacing the
        example data and code with your own.
  \item Compile with \texttt{markstat using yourfile, bundle} to produce
        \texttt{yourfile.html}.
\end{enumerate}
\end{tipbox}

\bigskip\horrule{0.4pt}\bigskip

%----------------------------------------------------------------------------------------
\section{What You Need to Install (One-Time Setup)}
%----------------------------------------------------------------------------------------

You need \textbf{four things} before you can compile your first Stata Markdown document.
Steps 1--3 are done inside Stata; Step 4 (Pandoc) is done outside Stata.

\medskip
\begin{center}
\begin{tabular}{clll}
\toprule
\textbf{Step} & \textbf{What} & \textbf{Where to install} & \textbf{Required?} \\
\midrule
1 & \texttt{markstat} Stata package & Inside Stata & Yes \\
2 & \texttt{whereis} Stata package  & Inside Stata & Yes \\
3 & Pandoc (external program)       & Your operating system & Yes \\
4 & Register Pandoc with Stata      & Inside Stata (once) & Yes \\
\bottomrule
\end{tabular}
\end{center}

\bigskip\horrule{0.4pt}\bigskip

%----------------------------------------------------------------------------------------
\section{Step 1 \& 2 --- Install Stata Packages}
%----------------------------------------------------------------------------------------

Open Stata and type the following two commands in the \textbf{Command window}:

\begin{lstlisting}[style=stata]
ssc install markstat, replace
ssc install whereis, replace
\end{lstlisting}

\begin{itemize}
  \item \texttt{markstat} is the package that reads your \texttt{.stmd} file, runs the
        Stata code blocks, and calls Pandoc to produce the final HTML.
  \item \texttt{whereis} is a helper that lets Stata find external programs (like Pandoc)
        on your computer.
  \item The \texttt{, replace} option safely re-installs even if the package already
        exists --- so it is always safe to run.
\end{itemize}

\begin{tipbox}
\textbf{Verify the installation:} After running the commands above, type
\texttt{which markstat} in the Command window. If it prints a file path, the package
was installed successfully.
\end{tipbox}

\bigskip\horrule{0.4pt}\bigskip

%----------------------------------------------------------------------------------------
\section{Step 3 --- Install Pandoc}
%----------------------------------------------------------------------------------------

\textbf{Pandoc} is a free, open-source document converter. It is what actually turns
your Stata output into a formatted HTML (or PDF) file. You must install it
\textit{outside} Stata, directly on your operating system.

\subsection*{3a. Download Pandoc}

Go to \url{https://pandoc.org/installing.html} and follow the instructions for your
operating system.

\medskip
\begin{warnbox}
\textbf{Important:} Pandoc is \textit{not} a Stata package --- it is a standalone
program. You install it just like any other application (e.g., Word, Chrome).
After installation, Stata does not automatically know where it is; you must
\textit{register} it manually (Step~4).
\end{warnbox}

\subsection*{3b. Windows Installation}

\begin{enumerate}
  \item Go to \url{https://pandoc.org/installing.html} and click the
        \textbf{Windows} installer link (\texttt{pandoc-x.x.x-windows-x86\_64.msi}).
  \item Run the downloaded \texttt{.msi} file and follow the installer wizard.
  \item By default, Pandoc installs to one of these locations:
\begin{lstlisting}[style=stata]
C:\Program Files\Pandoc\pandoc.exe
C:\Users\YourName\AppData\Local\Pandoc\pandoc.exe
\end{lstlisting}
  \item After installation, open \textbf{Command Prompt} (search ``cmd'' in the
        Start menu) and type:
\begin{lstlisting}[style=shell]
where pandoc
\end{lstlisting}
        This prints the exact path where Pandoc was installed.
        \textbf{Copy this path} --- you will need it in Step~4.
\end{enumerate}

\subsection*{3c. macOS Installation}

\textbf{Option A --- Homebrew (recommended):}
\begin{enumerate}
  \item If you do not have Homebrew, install it from \url{https://brew.sh}.
  \item Open \textbf{Terminal} and run:
\begin{lstlisting}[style=shell]
brew install pandoc
\end{lstlisting}
  \item Find the path:
\begin{lstlisting}[style=shell]
which pandoc
\end{lstlisting}
        Common results: \texttt{/usr/local/bin/pandoc} (Intel Mac) or
        \texttt{/opt/homebrew/bin/pandoc} (Apple Silicon Mac).
\end{enumerate}

\textbf{Option B --- Installer package:}
\begin{enumerate}
  \item Download the \texttt{.pkg} file from \url{https://pandoc.org/installing.html}.
  \item Double-click to run the installer.
  \item Find the path by opening Terminal and typing \texttt{which pandoc}.
\end{enumerate}

\subsection*{3d. Verify the Installation}

Regardless of your operating system, confirm Pandoc is installed by running in a
terminal / command prompt:

\begin{lstlisting}[style=shell]
pandoc --version
\end{lstlisting}

You should see something like \texttt{pandoc 3.x.x}. If you see an error, the
installation did not succeed --- repeat the steps above.

\bigskip\horrule{0.4pt}\bigskip

%----------------------------------------------------------------------------------------
\section{Step 4 --- Register Pandoc with Stata}
%----------------------------------------------------------------------------------------

This is the step most students miss. Even after installing Pandoc, Stata does not
automatically know where it is. You must tell Stata the full path using the
\texttt{whereis} command.

\medskip
\begin{warnbox}
\textbf{Run this command directly in the Stata Command window --- NOT inside a
\texttt{.do} file.} The registration is saved permanently on your machine, so you
only need to do this \textit{once} per computer.
\end{warnbox}

\subsection*{4a. Windows}

Using the path you found in Step 3b (adjust if yours differs):

\begin{lstlisting}[style=stata]
whereis pandoc "C:\Users\YourName\AppData\Local\Pandoc\pandoc.exe"
\end{lstlisting}

Or, if Pandoc was installed system-wide:

\begin{lstlisting}[style=stata]
whereis pandoc "C:\Program Files\Pandoc\pandoc.exe"
\end{lstlisting}

\subsection*{4b. macOS (Intel)}

\begin{lstlisting}[style=stata]
whereis pandoc "/usr/local/bin/pandoc"
\end{lstlisting}

\subsection*{4c. macOS (Apple Silicon)}

\begin{lstlisting}[style=stata]
whereis pandoc "/opt/homebrew/bin/pandoc"
\end{lstlisting}

\subsection*{4d. How to Verify}

After running \texttt{whereis pandoc}, type the same command again without the path
argument:

\begin{lstlisting}[style=stata]
whereis pandoc
\end{lstlisting}

If it prints the path you registered, you are good to go.

\bigskip\horrule{0.4pt}\bigskip

%----------------------------------------------------------------------------------------
\section{Step 5 --- Write Your \texttt{.stmd} File}
%----------------------------------------------------------------------------------------

A \texttt{.stmd} file is a plain text file. You can write it in any text editor
(Notepad, VS Code, Stata's Do-file editor, etc.) and save it with the
\texttt{.stmd} extension.

\subsection*{5a. Overall Structure}

Every \texttt{.stmd} file has three parts:

\begin{enumerate}
  \item \textbf{YAML header} --- document title, author, date (between \texttt{-{}-{}-} lines).
  \item \textbf{Markdown text} --- plain English narrative, formatted with Markdown syntax.
  \item \textbf{Stata code blocks} --- fenced with backticks, executed when compiled.
\end{enumerate}

\subsection*{5b. Minimal Example}

\begin{lstlisting}[style=stata]
---
title: "My First Stata Markdown"
author: "Your Name"
date: "2025"
---

## 1. Load Data

Here I load the dataset and display summary statistics.

```s
use "mydata.dta", clear
summarize
```

## 2. A Simple Graph

```s
histogram income, xtitle("Income") ytitle("Density")
graph export "income_hist.png", replace
```
\end{lstlisting}

\subsection*{5c. Types of Code Blocks}

\begin{center}
\begin{tabular}{lll}
\toprule
\textbf{Syntax} & \textbf{Code shown?} & \textbf{Output shown?} \\
\midrule
\texttt{```s}   & Yes & Yes \\
\texttt{```s/q} & No  & No  \quad (silent / quiet) \\
\texttt{```}    & Yes & Not executed \\
\bottomrule
\end{tabular}
\end{center}

\medskip
Use \texttt{```s/q} for setup code (path definitions, \texttt{ssc install}) that you
want to run silently without cluttering the report. Use \texttt{```s} for all
analysis steps whose output you want displayed.

\subsection*{5d. Path Setup (Portability Tip)}

Hard-coded file paths break when the file is opened on another computer. Use
\texttt{c(username)} to set paths conditionally:

\begin{lstlisting}[style=stata]
```s/q
if "`c(username)'" == "alice" {
    global rawdata = "C:\Users\alice\Desktop\data"
}
if "`c(username)'" == "bob" {
    global rawdata = "/Users/bob/Documents/data"
}
```
\end{lstlisting}

Then use \texttt{\$rawdata} throughout your code:

\begin{lstlisting}[style=stata]
```s
use "$rawdata/myfile.dta", clear
```
\end{lstlisting}

Find your username by typing \texttt{display "`c(username)'"} in the Stata Command
window.

\bigskip\horrule{0.4pt}\bigskip

%----------------------------------------------------------------------------------------
\section{Step 6 --- Compile the Document}
%----------------------------------------------------------------------------------------

Once your \texttt{.stmd} file is saved, compile it from the Stata Command window.

\subsection*{6a. Change to the Correct Directory}

Stata must be in the \textit{same folder} as your \texttt{.stmd} file when you compile:

\begin{lstlisting}[style=stata]
cd "C:\Users\YourName\Desktop\MyProject"
\end{lstlisting}

Or, if you defined \texttt{\$s\_markdown} in your path setup block:

\begin{lstlisting}[style=stata]
cd "$s_markdown"
\end{lstlisting}

\subsection*{6b. Compile to HTML}

\begin{lstlisting}[style=stata]
markstat using myfile, bundle
\end{lstlisting}

\begin{itemize}
  \item Replace \texttt{myfile} with the name of your \texttt{.stmd} file
        \textbf{without} the \texttt{.stmd} extension.
  \item The \texttt{bundle} option embeds all figures directly into the HTML file as
        Base64 data, producing a single self-contained file that you can share or
        submit without attaching any image files separately.
  \item Output: \texttt{myfile.html} in the same folder.
\end{itemize}

\subsection*{6c. Compile to PDF (Optional)}

\begin{lstlisting}[style=stata]
markstat using myfile, pdf
\end{lstlisting}

\begin{warnbox}
\textbf{PDF requires a \LaTeX{} distribution.} Install either
\href{https://miktex.org}{MiKTeX} (Windows) or
\href{https://tug.org/mactex/}{MacTeX} (macOS) before using the \texttt{pdf} option.
For homework submission, \textbf{HTML} (\texttt{bundle}) is simpler and recommended.
\end{warnbox}

\bigskip\horrule{0.4pt}\bigskip

%----------------------------------------------------------------------------------------
\section{Troubleshooting Common Errors}
%----------------------------------------------------------------------------------------

\begin{center}
\begin{tabular}{p{5.2cm} p{8.8cm}}
\toprule
\textbf{Error message} & \textbf{Cause and fix} \\
\midrule
\texttt{pandoc not found} &
Pandoc is not registered. Run \texttt{whereis pandoc "/full/path/to/pandoc"}
in the Stata Command window (see Step~4). \\[6pt]

\texttt{command markstat is unrecognized} &
\texttt{markstat} is not installed. Run \texttt{ssc install markstat, replace}. \\[6pt]

\texttt{command whereis is unrecognized} &
\texttt{whereis} is not installed. Run \texttt{ssc install whereis, replace}. \\[6pt]

\texttt{file not found} (at compile) &
Stata is not in the folder containing your \texttt{.stmd} file.
Run \texttt{cd "correct/path"} first. \\[6pt]

Image missing in HTML &
Use the \texttt{bundle} option: \texttt{markstat using myfile, bundle}.
Without it, images are external files that may not transfer. \\[6pt]

Variable already defined &
A code block ran twice. Add \texttt{capture drop varname} before \texttt{gen}. \\

\bottomrule
\end{tabular}
\end{center}

\bigskip\horrule{0.4pt}\bigskip

%----------------------------------------------------------------------------------------
\section{Quick Reference: Complete Setup Checklist}
%----------------------------------------------------------------------------------------

\begin{enumerate}
  \item[\(\square\)] \textbf{Install \texttt{markstat}} --- in Stata: \texttt{ssc install markstat, replace}
  \item[\(\square\)] \textbf{Install \texttt{whereis}} --- in Stata: \texttt{ssc install whereis, replace}
  \item[\(\square\)] \textbf{Install Pandoc} --- download from \url{https://pandoc.org/installing.html}
        (Windows installer or \texttt{brew install pandoc} on macOS)
  \item[\(\square\)] \textbf{Find the Pandoc path}
        \begin{itemize}
          \item Windows --- open Command Prompt, type \texttt{where pandoc}
          \item macOS --- open Terminal, type \texttt{which pandoc}
        \end{itemize}
  \item[\(\square\)] \textbf{Register Pandoc with Stata} (once per machine, in Command window):
\begin{lstlisting}[style=stata]
whereis pandoc "full\path\to\pandoc.exe"
\end{lstlisting}
  \item[\(\square\)] \textbf{Write} your \texttt{.stmd} file (YAML header + Markdown + \texttt{```s} blocks)
  \item[\(\square\)] \textbf{Compile} (in Stata Command window):
\begin{lstlisting}[style=stata]
cd "folder\containing\your\stmd\file"
markstat using yourfilename, bundle
\end{lstlisting}
  \item[\(\square\)] \textbf{Open} \texttt{yourfilename.html} in any web browser --- done!
\end{enumerate}

\bigskip\horrule{0.4pt}\bigskip

%----------------------------------------------------------------------------------------
\section{A Complete \texttt{.stmd} Template}
%----------------------------------------------------------------------------------------

Copy and save the following as \texttt{myreport.stmd}.
Edit the YAML header and add your own username block.

\begin{lstlisting}[style=stata]
---
title: "Your Report Title"
author: "Your Name"
date: "2025"
---

## 1. Setup

```s/q
clear all
set more off

* --- Set paths (add your username block here) ---
* Type  display "`c(username)'"  in Stata to find your username.

if "`c(username)'" == "alice" {
    global rawdata    = "C:\Users\alice\Desktop\data"
    global s_markdown = "C:\Users\alice\Desktop\data"
}
if "`c(username)'" == "bob" {
    global rawdata    = "/Users/bob/Documents/data"
    global s_markdown = "/Users/bob/Documents/data"
}
```

## 2. Load Data

```s
use "$rawdata/mydata.dta", clear
describe
summarize
```

## 3. Visualize

```s
histogram income, xtitle("Income") ytitle("Density")
graph export "$s_markdown/income_hist.png", replace
```

## 4. Regression

```s
reg outcome treatment control1 control2, r
```
\end{lstlisting}

\medskip
Compile with:

\begin{lstlisting}[style=stata]
cd "$s_markdown"
markstat using myreport, bundle
\end{lstlisting}

\end{document}
